\chapter{Large-\texorpdfstring{\(N\)}{N} Expansion and Gauge-Theory Strings}
\label{ch:volume-iii-large-n-expansion-gauge-theory-strings}

\section{The 't Hooft Limit}

Let the gauge group be \(SU(N)\), and write the Yang--Mills action in a
normalization where the interaction vertices are controlled by
\[
  \lambda=g_{\rm YM}^2N.
\]
The parameter \(\lambda\) is the 't Hooft coupling.  The 't Hooft limit is
\begin{equation}
  N\to\infty,
  \qquad
  \lambda\ \text{fixed}.
  \label{eq:t-hooft-limit}
\end{equation}
Adjoint fields are \(N\times N\) matrices.  Their propagators carry two color
indices, and it is useful to draw them as double lines.  A Feynman graph then
becomes a ribbon graph, hence a two-dimensional cell decomposition of an
oriented surface.

For a connected vacuum ribbon graph, let
\[
\begin{aligned}
  V&=\text{number of vertices},\\
  E&=\text{number of propagators},\\
  F&=\text{number of closed color-index faces}.
\end{aligned}
\]
The contribution scales as
\begin{equation}
  N^{F-E+V}\lambda^{E-V}
  =
  N^\chi\lambda^{E-V},
  \qquad
  \chi=V-E+F.
  \label{eq:ribbon-graph-euler-scaling}
\end{equation}
If the ribbon graph is drawn on a closed oriented surface of genus \(h\), then
\[
  \chi=2-2h.
\]
Thus planar graphs, \(h=0\), dominate; graphs of genus \(h\) are suppressed by
\(N^{-2h}\).

\begin{figure}[H]
  \centering
  \resizebox{0.94\textwidth}{!}{%
  \begin{tikzpicture}[x=1cm,y=1cm,every node/.style={font=\scriptsize}]
    \tikzset{arrow/.style={-{Latex[length=2mm]},line width=0.85pt}}
    \begin{scope}[shift={(-4.3,0)}]
      \node[font=\small] at (0,1.55) {ordinary graph};
      \draw[thick] (-1.0,0) .. controls (-0.5,0.75) and (0.5,0.75) .. (1.0,0);
      \draw[thick] (-1.0,0) .. controls (-0.5,-0.75) and (0.5,-0.75) .. (1.0,0);
      \draw[thick] (-0.32,0.55)--(-0.32,-0.55);
      \draw[thick] (0.32,0.55)--(0.32,-0.55);
    \end{scope}
    \node at (-2.35,0) {\(\longrightarrow\)};
    \begin{scope}[shift={(0,0)}]
      \node[font=\small] at (0,1.55) {double-line graph};
      \draw[thick] (-1.2,0) .. controls (-0.65,0.90) and (0.65,0.90) .. (1.2,0);
      \draw[thick] (-1.2,0) .. controls (-0.65,-0.90) and (0.65,-0.90) .. (1.2,0);
      \draw[thick] (-0.82,0.24) .. controls (-0.35,0.55) and (0.35,0.55) .. (0.82,0.24);
      \draw[thick] (-0.82,-0.24) .. controls (-0.35,-0.55) and (0.35,-0.55) .. (0.82,-0.24);
      \draw[blue!65!black,line width=0.9pt] (0,0) circle (0.42);
      \node[blue!65!black] at (0,-1.18) {index face};
    \end{scope}
    \draw[arrow] (1.75,0)--(2.85,0);
    \begin{scope}[shift={(4.45,0)}]
      \node[font=\small] at (0,1.55) {oriented surface};
      \draw[thick,fill=blue!6] (-1.25,-0.70) rectangle (1.25,0.70);
      \foreach \x in {-0.75,-0.25,0.25,0.75} {
        \draw[thick] (\x,-0.70)--(\x,0.70);
      }
      \node at (0,-1.18) {\(\chi=V-E+F\)};
    \end{scope}
  \end{tikzpicture}
  }
  \caption{Adjoint propagators carry two color lines.  The resulting ribbon
  graph has an Euler characteristic, and this characteristic controls the
  power of \(N\).}
  \label{fig:large-n-ribbon-graphs}
\end{figure}

\section{Single-Trace Normalization}

A single-trace operator of length \(L\) has the form
\[
  \operatorname{Tr}(X_1X_2\cdots X_L),
\]
where each \(X_\ell\) is an adjoint field or a covariant derivative of one.
It is useful to choose a normalization
\[
  \mathcal O_L
  =
  \frac1N
  \operatorname{Tr}(X_1X_2\cdots X_L)
\]
for large-\(N\) counting.  Connected correlators of \(k\) normalized
single-trace operators scale as
\begin{equation}
  \left\langle
    \mathcal O_{L_1}\cdots\mathcal O_{L_k}
  \right\rangle_{\rm conn}
  =
  O(N^{2-k})
  \label{eq:large-n-connected-correlator-scaling}
\end{equation}
at leading genus.  The next corrections are powers of \(N^{-2}\).  This is
the quantitative meaning of large-\(N\) factorization:
\[
  \langle \mathcal O_1\mathcal O_2\rangle
  =
  \langle \mathcal O_1\rangle\langle \mathcal O_2\rangle
  +O(N^{-2})
\]
for normalized single traces in a charge sector with nonzero one-point
functions.  When the one-point functions vanish, the connected two-point
function is the leading term and is \(O(N^0)\).

\section{Wilson Loops and Flux Strings}

Let \(C\) be an oriented closed curve in spacetime and \(R\) a representation
of \(G\).  The Wilson loop is
\begin{equation}
  W_R(C)
  =
  \operatorname{Tr}_R
  P\exp\left(\ii\oint_C A_\mu\,\dd x^\mu\right),
  \label{eq:large-n-wilson-loop}
\end{equation}
where \(P\) denotes path ordering along \(C\).  In a confining gauge theory,
large Wilson loops diagnose the flux string tension through an area law.
Large \(N\) suppresses splitting and joining of these flux strings, because
each splitting or joining changes the topology of the associated ribbon
surface and carries powers of \(1/N\).

For pure Yang--Mills, the running of the 't Hooft coupling generates a scale
\(\Lambda_{\rm YM}\).  Glueball masses and flux-string tensions are then
expected to be of the form
\[
  M_{\rm glueball}=c_M\Lambda_{\rm YM},
  \qquad
  T_{\rm flux}=c_T\Lambda_{\rm YM}^2,
\]
where \(c_M\) and \(c_T\) are dimensionless numbers determined by the theory.
Large \(N\) makes the interactions among color-singlet glueballs weak:
connected \(k\)-glueball amplitudes scale with decreasing powers of \(N\) as
\(k\) increases.

\section{Single Traces as Closed Chains}

In a matrix theory, cyclicity of the trace identifies
\[
  \operatorname{Tr}(X_1X_2\cdots X_L)
  =
  \operatorname{Tr}(X_2\cdots X_LX_1).
\]
Thus a length-\(L\) single-trace operator is naturally a closed chain of
\(L\) sites.  The field \(X_\ell\) placed at site \(\ell\) is the spin-chain
state at that site.  Planarity is the condition that, at leading order in
\(1/N\), local interactions in the trace act on neighboring sites.

\begin{figure}[H]
  \centering
  \resizebox{0.94\textwidth}{!}{%
  \begin{tikzpicture}[x=1cm,y=1cm,every node/.style={font=\scriptsize}]
    \tikzset{arrow/.style={-{Latex[length=2mm]},line width=0.85pt}}
    \begin{scope}[shift={(-4.2,0)}]
      \draw[thick] (0,0) ellipse (1.35 and 0.55);
      \foreach \ang/\lab in {20/{X_1},70/{X_2},125/{X_3},190/{X_4},250/{X_5},315/{X_L}} {
        \node[blue!65!black] at ({1.35*cos(\ang)},{0.55*sin(\ang)}) {\(\lab\)};
      }
      \node at (0,-1.25) {closed single trace};
    \end{scope}
    \draw[arrow] (-2.2,0)--(-1.0,0);
    \begin{scope}
      \draw[thick] (-0.45,0) -- (4.65,0);
      \foreach \x/\lab in {0/{X_1},0.85/{X_2},1.70/{X_3},2.55/{X_4},3.40/{X_5},4.25/{X_L}} {
        \draw[thick,fill=blue!5] (\x-0.28,-0.28) rectangle (\x+0.28,0.28);
        \node at (\x,0) {\(\lab\)};
      }
      \draw[orange!85!black,line width=0.9pt] (1.42,-0.48) rectangle (2.83,0.48);
      \node[orange!85!black] at (2.12,0.80) {planar local interaction};
      \node at (2.12,-1.25) {spin-chain representative};
    \end{scope}
  \end{tikzpicture}
  }
  \caption{A single trace is a closed chain because cyclic rotations of the
  fields define the same operator.  Planar interactions act locally on
  neighboring sites of this chain.}
  \label{fig:single-trace-closed-chain}
\end{figure}

This chain interpretation will become explicit in the planar dilatation
operator of \(\mathcal N=4\) super-Yang--Mills.
